\section{Prompt optimization must preserve the question, evidence and authority}
\label{sec:exec24-prompt-engineering}

An analyst can improve a prompt while making a banking answer less defensible.
More detailed instructions may introduce an unstated population assumption;
a supplied worked solution may make a benchmark easier without improving
independent problem solving; a tool-using agent may execute a valid command for
an invalid causal question. The relevant product question is therefore which
prompting design improves an explicitly defined service under a fixed evidence
and authority contract. Chen and colleagues' \emph{Unleashing the potential of
prompt engineering for large language models} provides a broad taxonomy of
instructions, demonstrations, decomposition, repeated sampling, search,
optimization, multimodal prompting, evaluation and security concerns.
\citep{exec24_promptengineering_original}

The inspected source is the 58-page arXiv revision 6 dated 11 May 2025, the
latest revision listed before the client's 7 October 2026 cutoff. The initial
23 October 2023 version has 25 pages and a different title; it is not silently
substituted for revision 6. The archive also supplies a journal identity, but
the final journal bytes have not been equated with this preprint. We use the
survey as a primary record of its authors' organization, examples and caveats.
Its references do not themselves qualify the methods, packages or empirical
rankings it describes. No model is run in this review.

\subsection{Extra context can change the target instead of identifying it}

The source's introductory examples compare a short question with more elaborate
prompts and generated context. In its entrance-width example, later answers
consider height, terrain, weight and access restrictions that were not fixed in
the question. Those may be useful questions to ask, but neither physical
possibility nor an expanded list of possibilities defines a probability law over
them. A preferred answer based on presumed prevalence imports an assumption.
The source describes richer answers as more coherent; our narrower conclusion
is that coherence must be evaluated separately from support by the given
information. Figure text was inspected, without fresh visual review.
\citep{exec24_promptengineering_original}

For a banking analogue, let $d$ be the reduction in default probability from an
intervention at a specified horizon, $L=1000$ dollars the loss per default and
$c=20$ dollars the action cost. Under the stipulated linear loss calculation,
net value is $V(d)=Ld-c$. Suppose the available information only restricts
$d\in[-0.10,0.10]$. The compatible values are
\[
 V(d)\in[-120,80]\quad\text{dollars per eligible account},
 \qquad V(d)>0\ \Longleftrightarrow\ d>0.02.
\]
If an elaborate prompt inserts $d=0.05$ as plausible generated knowledge, it
returns USD 30. The arithmetic is correct conditional on that value; the
supplied information did not identify it. The appropriate output retains the
range, break-even condition and the need for additional evidence. A user can
explicitly request the five-percentage-point scenario, but a scenario must be
labeled and must not overwrite the identified-set result.
\citep{exec24_promptengineering_reference}

The translation interface should therefore return missing-premise questions
alongside its candidate target and graph. It should distinguish a quoted fact,
an analyst-supplied restriction, an assumed scenario and model-generated
context. Retrieval adds a document, not automatic truth; the document's
population, time and authority still require interpretation. The proposal
prefers a structured questionnaire when that is sufficient. Generated knowledge
is suitable for proposing hypotheses and sensitivity cases, subject to explicit
status, rather than filling causal evidence gaps invisibly.

\subsection{A task graph is not a causal graph}

The survey discusses chains, trees and graphs of generated intermediate
information, and a decomposed workflow with a controller and subtask handlers.
These are descriptions of computation. An edge from a retrieved statement to
a proposed answer records a dependency of the answer on text; it does not
assert that intervening on an economic variable changes another economic
variable. The graph used to schedule computation and the structural graph used
for identification consequently need different types, identifiers and
validation rules. Converting one into the other requires a scientific argument.
\citep{exec24_promptengineering_original}

This distinction guides integration. A deterministic handler can parse a
versioned variable dictionary; an approved estimator can consume a reviewed
query; a model can propose an explanatory draft. Each handler should declare its
inputs, outputs, preconditions, failure statuses and artifact versions. The
controller must preserve an unidentified or refused status instead of translating
it into a numerical default. Repeating a subtask must create a traceable attempt,
not erase the earlier refusal. These are proposed CLI/MCP contracts, not a claim
that a native agent framework already enforces them.

An authored reliability example shows why local component scores do not settle
the end-to-end claim. Let $E_i$ mean that stage $i$ is correct, for five stages,
and stipulate $P(E_i)=19/20$ for every stage. If their correctness events are
independent, complete correctness has probability
\[
 P\!\left(\bigcap_{i=1}^{5}E_i\right)
     =\left(\frac{19}{20}\right)^5
     =\frac{2476099}{3200000}\simeq0.773781.
\]
With only those marginal probabilities known, the union bound and event
containment instead give
\[
 0.75\leq P\!\left(\bigcap_{i=1}^{5}E_i\right)\leq0.95.
\]
Disjoint five-percent failure events attain the lower bound; identical failure
events attain the upper bound. Independence is additional information, not a
default consequence of modular code. These are unconditional event assumptions
under a declared task law, not measurements of a five-stage application.
\citep{exec24_promptengineering_reference}

For sequential evaluation, later stages may receive wrong earlier outputs or
may never run after a refusal. Report both attempted-stage denominators and
complete-task outcomes under the real controller policy. In particular,
testing each handler only on correct supplied inputs establishes a different
property from testing the composed workflow on manager requests.

\subsection{Prompt selection is learning and needs held-out evaluation}

The survey covers selecting uncertain questions for annotation and optimizing
prompts using feedback, search and learned rewriting. It also reports that
different tasks, models and metrics limit comparisons, and that automated
prompting can fail to improve on simple baselines. These are source-level
observations, not reproduced rankings. An optimized prompt is fitted to
information: developer trials and expert edits are part of training even when
model weights stay fixed.\citep{exec24_promptengineering_original}

Consider a constructed null example with $K=20$ candidate prompts and $n=4$
development tasks. Every prompt has true success probability $1/2$. Assume
success indicators are independent across prompts and tasks. For prompt $j$,
let $S_j$ be its development successes. Then $S_j\sim\operatorname{Binomial}(4,1/2)$.
Selecting the largest score gives
\[
 P\{\max_j S_j=4\}=1-(15/16)^{20}\simeq0.724941.
\]
Writing $F(m)=P(S_j\leq m)$, the expected best development accuracy is
\[
 E\!\left[\frac{\max_j S_j}{4}\right]
 =\frac14\sum_{m=1}^{4}\{1-F(m-1)^{20}\}
 \simeq0.931096.
\]
Nevertheless, a fresh test independent of selection still has expected
accuracy $1/2$ for the selected prompt. A development winner above 93 percent
can arise in this deliberately small, heavily searched example without any
improvement in true performance. The numerical probabilities depend on the
declared independence law; they are not universal estimates of prompt overfitting.
\citep{exec24_promptengineering_reference}

The banking protocol should therefore reserve unseen tasks or task families,
record the complete search and repair budget, and separate development examples
from scoring evidence. Frozen model and retrieval versions make a comparison
interpretable but do not themselves guarantee generalization. Where tasks share
a template, customer history or policy source, that clustering affects the unit
of evaluation. Repeated paraphrases of one causal problem must not be counted
as independent scientific problems without justification.

The survey's discussion of supplied solution chains is particularly relevant:
it acknowledges that supplying ground-truth intermediate solutions limits what
a high solve rate demonstrates. Our evaluation should label an oracle-assisted
condition separately from unaided translation or identification. Likewise,
zero-shot describes the declared examples in a prompt, not verified absence of
the problem from a model's pretraining data. Neither a polished explanation nor
agreement among human judges independently establishes causal correctness.
\citep{exec24_promptengineering_original}

\subsection{Prompt text cannot grant tool authority}

The survey includes adversarial inputs, poisoned data, backdoors, prompt injection
and leakage. Some of its language is stronger than a bounded experiment can
warrant. This proposal retains the threat categories without adopting a claim
that prompt design ensures secure deployment. Full source inspection of an
attack or defense, its version and threat model would precede implementation.
\citep{exec24_promptengineering_original}

A proposed tool gate separates three logical conditions: $T$, a trusted request
authorizes the action; $A$, the action is allowed by the relevant policy; and
$H$, the request's artifact identities match those reviewed. Define admission
as $T\land A\land H$. Enumerating the eight Boolean configurations admits one
and rejects seven. This exact truth table demonstrates only the declared
predicate. It does not prove authentication, policy completeness, correct hash
binding, parser safety or atomic execution in a real system.
\citep{exec24_promptengineering_reference}

For example, retrieved text may recommend replacing a reviewed causal graph.
That text is evidence to inspect, not an authorized request to overwrite it.
A proposed MCP tool should return a candidate revision with source provenance;
an effect-producing action needs the separate trusted request and relevant
policy checks. The same distinction applies to CLI flags and shell commands.
Read-only retrieval, query construction and execution should have explicit
operation types, with the relevant authority verified at each external action.
Textual role instructions are useful guidance but are not the authority record.

\subsection{Conditional adoption and complete economic comparisons}

The choices are not simply an advanced prompt or no assistance. The proposal
compares a deterministic form, an expert-written template, retrieved examples,
bounded decomposition and adaptive optimization on the same declared tasks.
Use a simple template when it meets the target; add retrieval when source
coverage and provenance justify it; use original, qualified framework interfaces
when reusable orchestration earns its cost. A whole-framework port or compulsory
agent layer needs a demonstrated requirement. The survey supplies discovery
leads for these alternatives, not software rights or functioning integrations.

In one authored cost allocation, ten attempted tasks cost USD 4 in model calls,
USD 2 in retrieval and USD 24 in review. If three produce accepted outputs,
cost is USD 3 per attempt but USD 10 per accepted output. Actual pricing,
latency, maintenance, failures and correct-refusal value remain to be measured.
A cheap request is not automatically a cheap decision. Public documents and
constructed causal worlds permit this prototype comparison before authentic
bank experiments, while source rights, representation and transfer limits remain
explicit.\citep{exec24_promptengineering_reference}

The resulting adoption rule is conditional: keep a prompt or agent design only
when its declared target, evidence handling, complete-task behavior and cost
justify it against the simpler alternative. This source broadens the taxonomy
and sharpens test requirements; it leaves native integration, complete frontier
coverage and bank acceptance unestablished.
